\vssub
\subsubsection{~三阶格式 (UQ)}
\opthead{UQ}{\ws}{H. L. Tolman}

\noindent
$\theta$ 空间中的 \uq{} 格式采用与物理空间格式类似的方式实现，区别在于封闭的方向空间不需要边界条件。$k$ 空间中的可变网格间距需要对格式作若干修改，如 \cite[{Appendix}]{art:Leo79} 所述。于是式~(\ref{eq:quick_1}) 至
(\ref{eq:quick_4}) 变为

% ------ QUICKEST scheme for k space---------- %
% eq:quick_1k        Basic flux
% eq:quick_2k        Boundary value
% eq:quick_3k        Divergence
% eq:quick_4k        CFL number

\begin{equation}
\cF_{m,-} = \left [ \dot{k}_{g,b} \: N_b \: \right ]^n_{i,j,l}
\: , \label{eq:quick_1k}\end{equation} \begin{equation}
\dot{k}_{g,b} = 0.5 \: \left ( \dot{k}_{g,m-1} + \dot{k}_{g,m} 
\: \right )  \: , \label{eq:quick_1ak}
\end{equation} \begin{equation}
N_b = \frac{1}{2} \left [ \rule[0mm]{0mm}{\baselineskip} \: 
(1+C)N_{i-1} + (1-C)N_i \: \right ] - \:
\frac{1-C^2}{6} \: {\cal CU} \: \Delta k^2_{m-1/2}, \label{eq:quick_2k} \end{equation} \begin{equation}
{\cal CU} =  \left \{ \begin{array}{ccc}
\frac{1}{\Delta k_{m-1}}
\left [ \frac{N_{ m }-N_{m-1}}{\Delta k_{m-1/2}} - 
        \frac{N_{m-1}-N_{m-2}}{\Delta k_{m,-3/2}} \right ]
               & \mbox{for} & \dot{k}_b \geq 0 \\
\frac{1}{\Delta k_m}
\left [ \frac{N_{m+1}-N_{ m }}{\Delta k_{m+1/2}} -
       \frac{N_{ m }-N_{m-1}}{\Delta k_{m-1/2}} \right ]
               & \mbox{for} & \dot{k}_b   <  0
\end{array} \right . \: , \label{eq:quick_3k}
\end{equation} \begin{equation}
C = \frac{\dot{k}_{g,b} \: \Delta t}{\Delta k_{m-1/2}}
\: , \label{eq:quick_4k} \end{equation}

\noindent
其中 $\Delta k_m$ 是网格点 $m$ 处的离散谱带或单元宽度，$\Delta k_{m-1/2}$ 是计数器为 $m$ 和 $m-1$ 的网格点之间的距离。若使用式~(\ref{eq:quick_4k}) 的 \cfl{} 数，则可像式~(\ref{eq:ult_1}) 至
(\ref{eq:ult_4}) 那样应用 \ult{} 限制器。在低波数和高波数边界处，通量仍用一阶迎风方法估计，边界条件采用上文为一阶格式定义的条件。$k$ 空间中的最终格式为

% eq:1uq_k_tot

\begin{equation}
N_{i,j,l,m}^{n+1} = N_{i,j,l,m}^n 
 + \frac{\Delta t}{\Delta k_m} \left [ \cF_{m,-} - \cF_{m,+} \right ]
\: , \label{eq:uq_k_tot} \end{equation}
